
% \newpage
This section presents the Pourbaix diagrams generated in this work. For a given metal, ligand pair, if no formation energies or formation constants were available for aqueous metal, ligand complexes, that specific ligand was excluded by setting its concentration to zero in the corresponding diagram. Other ligands with available complexation data were still included where applicable.

%%%%%%%%%%%%%%%%%%%%%%%%%%%%%%%%% Au %%%%%%%%%%%%%%%%%%%%%%%%%%%%%%%%%
\begin{figure}[!htbp]
    \centering
    \includegraphics[width=\textwidth]{SI_Au_3panel_side_colors.pdf}
    \caption{Au Pourbaix diagrams. The metal-ion activity is $10^{-4}\,\mathrm{M}$. (a) No added ligands; (b) \ce{NH3} and Gly; (c) \ce{NH3}, Gly, and \ce{CN-}. The ligand-containing panels use $[\ce{NH3}]_{\mathrm{tot}}=0.02\,\mathrm{M}$ and $[\mathrm{Gly}]_{\mathrm{tot}}=0.1\,\mathrm{M}$; panel (c) also uses $[\ce{CN-}]_{\mathrm{tot}}=10^{-4}\,\mathrm{M}$. The black outline marks the experimental window, pH 10.5--13.5 and $-2.0$ to $2.3$~V vs RHE (converted to SHE on the plotted potential axis).}
    \label{fig:Au_Pourbaix_SI}
    % Retain existing panel cross-references for the combined artwork.
    \phantomsubcaption\label{fig:Au_Pourbaix_H2O}
    \phantomsubcaption\label{fig:Au_Pourbaix_NH3_Gly_SI}
    \phantomsubcaption\label{fig:Au_Pourbaix_NH3_Gly_CN_SI}
\end{figure}

%%%%%%%%%%%%%%%%%%%%%%%%%%%%%%%%% Pd %%%%%%%%%%%%%%%%%%%%%%%%%%%%%%%%%
\begin{figure}[!htbp]
    \centering
    \includegraphics[width=\textwidth]{SI_Pd_3panel_side_colors.pdf}
    \caption{Pd Pourbaix diagrams. The metal-ion activity is $10^{-4}\,\mathrm{M}$. (a) No added ligands; (b,c) \ce{NH3}, Gly, and \ce{CN-}. Panels (b,c) use $[\ce{NH3}]_{\mathrm{tot}}=0.02\,\mathrm{M}$, $[\mathrm{Gly}]_{\mathrm{tot}}=0.1\,\mathrm{M}$, and $[\ce{CN-}]_{\mathrm{tot}}=10^{-4}\,\mathrm{M}$. Panel (b) uses the historical \ce{Pd(CN)4^2-} energy; panel (c) uses the tabulated energy based on $\log\beta_4=62.3$ reported by Harrington et al.~\cite{Harrington2005DeterminationIon}. The black outline marks the experimental window, pH 10.5--13.5 and $-2.0$ to $2.3$~V vs RHE (converted to SHE on the plotted potential axis).}
    \label{fig:Pd_Pourbaix_SI}
    % Retain existing panel cross-references for the combined artwork.
    \phantomsubcaption\label{fig:Pd_Pourbaix_NH3_Gly_SI}
    \phantomsubcaption\label{fig:Pd_Pourbaix_NH3_Gly_CN_SI}
    \phantomsubcaption\label{fig:Pd_Pourbaix_NH3_Gly_CN_SI_Harrington}
\end{figure}

%%%%%%%%%%%%%%%%%%%%%%%%%%%%%%%%% Ni %%%%%%%%%%%%%%%%%%%%%%%%%%%%%%%%%
\begin{figure}[!htbp]
    \centering
    \includegraphics[width=\textwidth]{SI_Ni_3panel_side_colors.pdf}
    \caption{Ni Pourbaix diagrams. The metal-ion activity is $10^{-4}\,\mathrm{M}$. (a) No added ligands; (b) \ce{NH3} and Gly; (c) \ce{NH3}, Gly, and \ce{CN-}. The ligand-containing panels use $[\ce{NH3}]_{\mathrm{tot}}=0.02\,\mathrm{M}$ and $[\mathrm{Gly}]_{\mathrm{tot}}=0.1\,\mathrm{M}$; panel (c) also uses $[\ce{CN-}]_{\mathrm{tot}}=10^{-4}\,\mathrm{M}$. The black outline marks the experimental window, pH 10.5--13.5 and $-2.0$ to $2.3$~V vs RHE (converted to SHE on the plotted potential axis).}
    \label{fig:Ni_Pourbaix_SI}
    % Retain existing panel cross-references for the combined artwork.
    \phantomsubcaption\label{fig:Ni_Pourbaix_H2O}
    \phantomsubcaption\label{fig:Ni_Pourbaix_NH3_Gly_SI}
    \phantomsubcaption\label{fig:Ni_Pourbaix_NH3_Gly_CN_SI}
\end{figure}

%%%%%%%%%%%%%%%%%%%%%%%%%%%%%%%%% Cu %%%%%%%%%%%%%%%%%%%%%%%%%%%%%%%%%
\begin{figure}[!htbp]
    \centering
    \includegraphics[width=\textwidth]{SI_Cu_3panel_side_colors.pdf}
    \caption{Cu Pourbaix diagrams. The metal-ion activity is $10^{-4}\,\mathrm{M}$. (a) No added ligands; (b) \ce{NH3} and Gly; (c) \ce{NH3}, Gly, and \ce{CN-}. The ligand-containing panels use $[\ce{NH3}]_{\mathrm{tot}}=0.02\,\mathrm{M}$ and $[\mathrm{Gly}]_{\mathrm{tot}}=0.1\,\mathrm{M}$; panel (c) also uses $[\ce{CN-}]_{\mathrm{tot}}=10^{-4}\,\mathrm{M}$. The black outline marks the experimental window, pH 10.5--13.5 and $-2.0$ to $2.3$~V vs RHE (converted to SHE on the plotted potential axis).}
    \label{fig:Cu_Pourbaix_SI}
    % Retain existing panel cross-references for the combined artwork.
    \phantomsubcaption\label{fig:Cu_Pourbaix_H2O}
    \phantomsubcaption\label{fig:Cu_Pourbaix_NH3_Gly_SI}
    \phantomsubcaption\label{fig:Cu_Pourbaix_NH3_Gly_CN_SI}
\end{figure}

%%%%%%%%%%%%%%%%%%%%%%%%%%%%%%%%% Fe %%%%%%%%%%%%%%%%%%%%%%%%%%%%%%%%%
\begin{figure}[!htbp]
    \centering
    \includegraphics[width=\textwidth]{SI_Fe_3panel_side_colors.pdf}
    \caption{Fe Pourbaix diagrams. The metal-ion activity is $10^{-4}\,\mathrm{M}$. (a) No added ligands; (b) \ce{NH3} and Gly; (c) \ce{NH3}, Gly, and \ce{CN-}. The ligand-containing panels use $[\ce{NH3}]_{\mathrm{tot}}=0.02\,\mathrm{M}$ and $[\mathrm{Gly}]_{\mathrm{tot}}=0.1\,\mathrm{M}$; panel (c) also uses $[\ce{CN-}]_{\mathrm{tot}}=10^{-4}\,\mathrm{M}$. The black outline marks the experimental window, pH 10.5--13.5 and $-2.0$ to $2.3$~V vs RHE (converted to SHE on the plotted potential axis).}
    \label{fig:Fe_Pourbaix_SI}
    % Retain existing panel cross-references for the combined artwork.
    \phantomsubcaption\label{fig:Fe_Pourbaix_H2O}
    \phantomsubcaption\label{fig:Fe_Pourbaix_NH3_Gly_SI}
    \phantomsubcaption\label{fig:Fe_Pourbaix_NH3_Gly_CN_SI}
\end{figure}

%%%%%%%%%%%%%%%%%%%%%%%%%%%%%%%%% Mn %%%%%%%%%%%%%%%%%%%%%%%%%%%%%%%%%
\begin{figure}[!htbp]
    \centering
    \includegraphics[width=\textwidth]{SI_Mn_2panel_side_colors.pdf}
    \caption{Mn Pourbaix diagrams. The metal-ion activity is $10^{-4}\,\mathrm{M}$. (a) No added ligands; (b) \ce{NH3} and Gly. The ligand-containing panels use $[\ce{NH3}]_{\mathrm{tot}}=0.02\,\mathrm{M}$ and $[\mathrm{Gly}]_{\mathrm{tot}}=0.1\,\mathrm{M}$. The black outline marks the experimental window, pH 10.5--13.5 and $-2.0$ to $2.3$~V vs RHE (converted to SHE on the plotted potential axis).}
    \label{fig:Mn_Pourbaix_SI}
    % Retain existing panel cross-references for the combined artwork.
    \phantomsubcaption\label{fig:Mn_Pourbaix_H2O}
    \phantomsubcaption\label{fig:Mn_Pourbaix_NH3_Gly_SI}
\end{figure}

%%%%%%%%%%%%%%%%%%%%%%%%%%%%%%%%% Zn %%%%%%%%%%%%%%%%%%%%%%%%%%%%%%%%%
\begin{figure}[!htbp]
    \centering
    \includegraphics[width=\textwidth]{SI_Zn_3panel_side_colors.pdf}
    \caption{Zn Pourbaix diagrams. The metal-ion activity is $10^{-4}\,\mathrm{M}$. (a) No added ligands; (b) \ce{NH3} and Gly; (c) \ce{NH3}, Gly, and \ce{CN-}. The ligand-containing panels use $[\ce{NH3}]_{\mathrm{tot}}=0.02\,\mathrm{M}$ and $[\mathrm{Gly}]_{\mathrm{tot}}=0.1\,\mathrm{M}$; panel (c) also uses $[\ce{CN-}]_{\mathrm{tot}}=10^{-4}\,\mathrm{M}$. The black outline marks the experimental window, pH 10.5--13.5 and $-2.0$ to $2.3$~V vs RHE (converted to SHE on the plotted potential axis).}
    \label{fig:Zn_Pourbaix_SI}
    % Retain existing panel cross-references for the combined artwork.
    \phantomsubcaption\label{fig:Zn_Pourbaix_H2O}
    \phantomsubcaption\label{fig:Zn_Pourbaix_NH3_Gly_SI}
    \phantomsubcaption\label{fig:Zn_Pourbaix_NH3_Gly_CN_SI}
\end{figure}

%%%%%%%%%%%%%%%%%%%%%%%%%%%%%%%%% Cd %%%%%%%%%%%%%%%%%%%%%%%%%%%%%%%%%
\begin{figure}[!htbp]
    \centering
    \includegraphics[width=\textwidth]{SI_Cd_3panel_side_colors.pdf}
    \caption{Cd Pourbaix diagrams. The metal-ion activity is $10^{-4}\,\mathrm{M}$. (a) No added ligands; (b) \ce{NH3} and Gly; (c) \ce{NH3}, Gly, and \ce{CN-}. The ligand-containing panels use $[\ce{NH3}]_{\mathrm{tot}}=0.02\,\mathrm{M}$ and $[\mathrm{Gly}]_{\mathrm{tot}}=0.1\,\mathrm{M}$; panel (c) also uses $[\ce{CN-}]_{\mathrm{tot}}=10^{-4}\,\mathrm{M}$. The black outline marks the experimental window, pH 10.5--13.5 and $-2.0$ to $2.3$~V vs RHE (converted to SHE on the plotted potential axis).}
    \label{fig:Cd_Pourbaix_SI}
    % Retain existing panel cross-references for the combined artwork.
    \phantomsubcaption\label{fig:Cd_Pourbaix_H2O}
    \phantomsubcaption\label{fig:Cd_Pourbaix_NH3_Gly_SI}
    \phantomsubcaption\label{fig:Cd_Pourbaix_NH3_Gly_CN_SI}
\end{figure}

%%%%%%%%%%%%%%%%%%%%%%%%%%%%%%%%% Co %%%%%%%%%%%%%%%%%%%%%%%%%%%%%%%%%
\begin{figure}[!htbp]
    \centering
    \includegraphics[width=\textwidth]{SI_Co_3panel_side_colors.pdf}
    \caption{Co Pourbaix diagrams. The metal-ion activity is $10^{-4}\,\mathrm{M}$. (a) No added ligands; (b) \ce{NH3} and Gly; (c) \ce{NH3}, Gly, and \ce{CN-}. The ligand-containing panels use $[\ce{NH3}]_{\mathrm{tot}}=0.02\,\mathrm{M}$ and $[\mathrm{Gly}]_{\mathrm{tot}}=0.1\,\mathrm{M}$; panel (c) also uses $[\ce{CN-}]_{\mathrm{tot}}=10^{-4}\,\mathrm{M}$. The black outline marks the experimental window, pH 10.5--13.5 and $-2.0$ to $2.3$~V vs RHE (converted to SHE on the plotted potential axis).}
    \label{fig:Co_Pourbaix_SI}
    % Retain existing panel cross-references for the combined artwork.
    \phantomsubcaption\label{fig:Co_Pourbaix_H2O}
    \phantomsubcaption\label{fig:Co_Pourbaix_NH3_Gly_SI}
    \phantomsubcaption\label{fig:Co_Pourbaix_NH3_Gly_CN_SI}
\end{figure}

%%%%%%%%%%%%%%%%%%%%%%%%%%%%%%%%% Sr %%%%%%%%%%%%%%%%%%%%%%%%%%%%%%%%%
\begin{figure}[!htbp]
    \centering
    \includegraphics[width=0.6\textwidth]{SI_Sr_1panel_side_colors.pdf}
    \caption{Sr Pourbaix diagram at $[\ce{NH3}]_{\mathrm{tot}}=0.02\,\mathrm{M}$ and $[\mathrm{Gly}]_{\mathrm{tot}}=0.1\,\mathrm{M}$, with no added cyanide. The metal-ion activity is $10^{-4}\,\mathrm{M}$. \ce{[Sr(Gly)]^2+} is the only included Sr--glycine complex because only one corresponding formation constant was available. The black outline marks the experimental window, pH 10.5--13.5 and $-2.0$ to $2.3$~V vs RHE (converted to SHE on the plotted potential axis).}
    \label{fig:Sr_Pourbaix_Harrington_SI}
\end{figure}
%%%%%%%%%%%%%%%%%%%%%%%%%%%%%%%% Pt Harrington %%%%%%%%%%%%%%%%%%%%%%%%%%%%%%%%%

% \begin{figure}[!htbp]
%     \centering
%     \includegraphics[width=0.7\textwidth]{Figures/pourbaix_diagrams/Pt/Pt-NH3-H2O_T=298.15_activity=1e-04_[NH3]=0.02M_[Gly]=0.1M_[CN]=0.0001_Harrington.png}
%     \caption{Pourbaix diagram for Pt at $[NH_3]_{initial}= 0.02M$, $[Gly]_{initial}=0.1M$ and $[\ce{CN-}]_{initial}=1e^{-4}$M. \ce{[Pt(CN)4]^2-} formation energy calculated using estimated stability constant $log\beta_4=70$ by Harring et al.\cite{Harrington2005DeterminationIon} Black box indicates experimental condition at applied potential vs RHE = -2 to 2 V, pH = 11.5 to 13.5.}
%     \label{fig:Pt_Pourbaix_Harrington_SI}
% \end{figure}
%%%%%%%%%%%%%%%%%%%%%%%%%%%%%% Ti alloys %%%%%%%%%%%%%%%%%%%%%%%%%%%%%%
\begin{figure}[!htbp]
    \centering
    \includegraphics[width=\textwidth]{SI_TiNi_2panel_side_colors_all_species.pdf}
    \caption{Pourbaix diagrams of TiNi. (a) \ce{NH3} and Gly; (b) \ce{NH3}, Gly, and \ce{CN-}. The metal-ion activity is $10^{-4}$. Ligand-containing panels use $[\ce{NH3}]_{\mathrm{tot}}=0.02\,\mathrm{M}$ and $[\mathrm{Gly}]_{\mathrm{tot}}=0.1\,\mathrm{M}$; cyanide-containing panels additionally use $[\ce{CN-}]_{\mathrm{tot}}=10^{-4}\,\mathrm{M}$. The black outline marks pH 10.5--13.5 and $-2.0$ to $2.3$ V vs RHE (converted to SHE and clipped to the plotted range). The side legends list all stable species combinations, including solids, aqueous ions, and metal--ligand complexes. The dashed blue and red lines mark the water reduction and oxidation boundaries, respectively.}
    \label{fig:Ti_Ni_Pourbaix_SI}
    \phantomsubcaption\label{fig:TiNi_Pourbaix_NH3_Gly_SI}
    \phantomsubcaption\label{fig:TiNi_Pourbaix_NH3_Gly_CN_SI}
\end{figure}

\begin{figure}[!htbp]
    \centering
    \includegraphics[width=\textwidth]{SI_TiCu_2panel_side_colors_all_species.pdf}
    \caption{Pourbaix diagrams of TiCu. (a) No added ligands; (b) \ce{NH3}, Gly, and \ce{CN-}. The metal-ion activity is $10^{-4}$. Ligand-containing panels use $[\ce{NH3}]_{\mathrm{tot}}=0.02\,\mathrm{M}$ and $[\mathrm{Gly}]_{\mathrm{tot}}=0.1\,\mathrm{M}$; cyanide-containing panels additionally use $[\ce{CN-}]_{\mathrm{tot}}=10^{-4}\,\mathrm{M}$. The black outline marks pH 10.5--13.5 and $-2.0$ to $2.3$ V vs RHE (converted to SHE and clipped to the plotted range). The side legends list all stable species combinations, including solids, aqueous ions, and metal--ligand complexes. The dashed blue and red lines mark the water reduction and oxidation boundaries, respectively.}
    \label{fig:Ti_Cu_Pourbaix_SI}
    \phantomsubcaption\label{fig:TiCu_Pourbaix_NH3_Gly_SI}
    \phantomsubcaption\label{fig:TiCu_Pourbaix_NH3_Gly_CN_SI}
\end{figure}
% Ti-Fe
\begin{figure}[!htbp]
    \centering
    \includegraphics[width=\textwidth]{SI_TiFe_2panel_side_colors_all_species.pdf}
    \caption{Pourbaix diagrams of TiFe. (a) \ce{NH3} and Gly; (b) \ce{NH3}, Gly, and \ce{CN-}. The metal-ion activity is $10^{-4}$. Ligand-containing panels use $[\ce{NH3}]_{\mathrm{tot}}=0.02\,\mathrm{M}$ and $[\mathrm{Gly}]_{\mathrm{tot}}=0.1\,\mathrm{M}$; cyanide-containing panels additionally use $[\ce{CN-}]_{\mathrm{tot}}=10^{-4}\,\mathrm{M}$. The black outline marks pH 10.5--13.5 and $-2.0$ to $2.3$ V vs RHE (converted to SHE and clipped to the plotted range). The side legends list all stable species combinations, including solids, aqueous ions, and metal--ligand complexes. The dashed blue and red lines mark the water reduction and oxidation boundaries, respectively.}
    \label{fig:Ti_Fe_Pourbaix_SI}
    \phantomsubcaption\label{fig:TiFe_Pourbaix_NH3_Gly_SI}
    \phantomsubcaption\label{fig:TiFe_Pourbaix_NH3_Gly_CN_SI}
\end{figure}

\begin{figure}[!htbp]
    \centering
    \includegraphics[width=0.6\textwidth]{SI_TiMg_1panel_side_colors_all_species.pdf}
    \caption{Pourbaix diagram of TiMg. (a) \ce{NH3} and Gly. The metal-ion activity is $10^{-4}$. Ligand-containing panels use $[\ce{NH3}]_{\mathrm{tot}}=0.02\,\mathrm{M}$ and $[\mathrm{Gly}]_{\mathrm{tot}}=0.1\,\mathrm{M}$, with no added cyanide. The black outline marks pH 10.5--13.5 and $-2.0$ to $2.3$ V vs RHE (converted to SHE and clipped to the plotted range). The side legends list all stable species combinations, including solids, aqueous ions, and metal--ligand complexes. The dashed blue and red lines mark the water reduction and oxidation boundaries, respectively.}
    \label{fig:Ti_Mg_Pourbaix_SI}
    \phantomsubcaption\label{fig:TiMg_Pourbaix_NH3_Gly_SI}
\end{figure}

% Ti-Mn
\begin{figure}[!htbp]
    \centering
    \includegraphics[width=0.6\textwidth]{SI_TiMn_1panel_side_colors_all_species.pdf}
    \caption{Pourbaix diagram of TiMn. (a) \ce{NH3} and Gly. The metal-ion activity is $10^{-4}$. Ligand-containing panels use $[\ce{NH3}]_{\mathrm{tot}}=0.02\,\mathrm{M}$ and $[\mathrm{Gly}]_{\mathrm{tot}}=0.1\,\mathrm{M}$, with no added cyanide. The black outline marks pH 10.5--13.5 and $-2.0$ to $2.3$ V vs RHE (converted to SHE and clipped to the plotted range). The side legends list all stable species combinations, including solids, aqueous ions, and metal--ligand complexes. The dashed blue and red lines mark the water reduction and oxidation boundaries, respectively.}
    \label{fig:Ti_Mn_Pourbaix_SI}
    \phantomsubcaption\label{fig:TiMn_Pourbaix_NH3_Gly_SI}
\end{figure}


% Ti-Zn
\begin{figure}[!htbp]
    \centering
    \includegraphics[width=0.6\textwidth]{SI_TiZn_1panel_side_colors_all_species.pdf}
    \caption{Pourbaix diagram of TiZn. (a) \ce{NH3}, Gly, and \ce{CN-}. The metal-ion activity is $10^{-4}$. Ligand-containing panels use $[\ce{NH3}]_{\mathrm{tot}}=0.02\,\mathrm{M}$ and $[\mathrm{Gly}]_{\mathrm{tot}}=0.1\,\mathrm{M}$; cyanide-containing panels additionally use $[\ce{CN-}]_{\mathrm{tot}}=10^{-4}\,\mathrm{M}$. The black outline marks pH 10.5--13.5 and $-2.0$ to $2.3$ V vs RHE (converted to SHE and clipped to the plotted range). The side legends list all stable species combinations, including solids, aqueous ions, and metal--ligand complexes. The dashed blue and red lines mark the water reduction and oxidation boundaries, respectively.}
    \label{fig:Ti_Zn_Pourbaix_SI}
    \phantomsubcaption\label{fig:TiZn_Pourbaix_NH3_Gly_CN_SI}
\end{figure}




%%%%%%%%%%%%%%%%%%%%%%%%%%%%%%%% Pd Alloys %%%%%%%%%%%%%%%%%%%%%%%%%%%%%%%%%
\begin{figure}[!htbp]
    \centering
    \includegraphics[width=\textwidth]{SI_PdAg_2panel_side_colors_all_species.pdf}
    \caption{Pourbaix diagrams of PdAg. (a) \ce{NH3} and Gly; (b) \ce{NH3}, Gly, and \ce{CN-}. The metal-ion activity is $10^{-4}$. Ligand-containing panels use $[\ce{NH3}]_{\mathrm{tot}}=0.02\,\mathrm{M}$ and $[\mathrm{Gly}]_{\mathrm{tot}}=0.1\,\mathrm{M}$; cyanide-containing panels additionally use $[\ce{CN-}]_{\mathrm{tot}}=10^{-4}\,\mathrm{M}$. The black outline marks pH 10.5--13.5 and $-2.0$ to $2.3$ V vs RHE (converted to SHE and clipped to the plotted range). The side legends list all stable species combinations, including solids, aqueous ions, and metal--ligand complexes. The dashed blue and red lines mark the water reduction and oxidation boundaries, respectively.}
    \label{fig:Pd_Ag_Pourbaix_SI}
    \phantomsubcaption\label{fig:PdAg_Pourbaix_NH3_Gly_SI}
    \phantomsubcaption\label{fig:PdAg_Pourbaix_NH3_Gly_CN_SI}
\end{figure}

\begin{figure}[!htbp]
    \centering
    \includegraphics[width=\textwidth]{SI_PdCu_2panel_side_colors_all_species.pdf}
    \caption{Pourbaix diagrams of PdCu. (a) \ce{NH3} and Gly; (b) \ce{NH3}, Gly, and \ce{CN-}. The metal-ion activity is $10^{-4}$. Ligand-containing panels use $[\ce{NH3}]_{\mathrm{tot}}=0.02\,\mathrm{M}$ and $[\mathrm{Gly}]_{\mathrm{tot}}=0.1\,\mathrm{M}$; cyanide-containing panels additionally use $[\ce{CN-}]_{\mathrm{tot}}=10^{-4}\,\mathrm{M}$. The black outline marks pH 10.5--13.5 and $-2.0$ to $2.3$ V vs RHE (converted to SHE and clipped to the plotted range). The side legends list all stable species combinations, including solids, aqueous ions, and metal--ligand complexes. The dashed blue and red lines mark the water reduction and oxidation boundaries, respectively.}
    \label{fig:Pd_Cu_Pourbaix_SI}
    \phantomsubcaption\label{fig:PdCu_Pourbaix_NH3_Gly_SI}
    \phantomsubcaption\label{fig:PdCu_Pourbaix_NH3_Gly_CN_SI}
\end{figure}

\begin{figure}[!htbp]
    \centering
    \includegraphics[width=\textwidth]{SI_PdNi_2panel_side_colors_all_species.pdf}
    \caption{Pourbaix diagrams of PdNi. (a) \ce{NH3} and Gly; (b) \ce{NH3}, Gly, and \ce{CN-}. The metal-ion activity is $10^{-4}$. Ligand-containing panels use $[\ce{NH3}]_{\mathrm{tot}}=0.02\,\mathrm{M}$ and $[\mathrm{Gly}]_{\mathrm{tot}}=0.1\,\mathrm{M}$; cyanide-containing panels additionally use $[\ce{CN-}]_{\mathrm{tot}}=10^{-4}\,\mathrm{M}$. The black outline marks pH 10.5--13.5 and $-2.0$ to $2.3$ V vs RHE (converted to SHE and clipped to the plotted range). The side legends list all stable species combinations, including solids, aqueous ions, and metal--ligand complexes. The dashed blue and red lines mark the water reduction and oxidation boundaries, respectively.}
    \label{fig:Pd_Ni_Pourbaix_SI}
    \phantomsubcaption\label{fig:PdNi_Pourbaix_NH3_Gly_SI}
    \phantomsubcaption\label{fig:PdNi_Pourbaix_NH3_Gly_CN_SI}
\end{figure}

\begin{figure}[!htbp]
    \centering
    \includegraphics[width=\textwidth]{SI_PdPt_2panel_side_colors_all_species.pdf}
    \caption{Pourbaix diagrams of PdPt. (a) \ce{NH3} and Gly; (b) \ce{NH3}, Gly, and \ce{CN-}. The metal-ion activity is $10^{-4}$. Ligand-containing panels use $[\ce{NH3}]_{\mathrm{tot}}=0.02\,\mathrm{M}$ and $[\mathrm{Gly}]_{\mathrm{tot}}=0.1\,\mathrm{M}$; cyanide-containing panels additionally use $[\ce{CN-}]_{\mathrm{tot}}=10^{-4}\,\mathrm{M}$. The black outline marks pH 10.5--13.5 and $-2.0$ to $2.3$ V vs RHE (converted to SHE and clipped to the plotted range). The side legends list all stable species combinations, including solids, aqueous ions, and metal--ligand complexes. The dashed blue and red lines mark the water reduction and oxidation boundaries, respectively.}
    \label{fig:Pd_Pt_Pourbaix_SI}
    \phantomsubcaption\label{fig:PdPt_Pourbaix_NH3_Gly_SI}
    \phantomsubcaption\label{fig:PdPt_Pourbaix_NH3_Gly_CN_SI}
\end{figure}

\begin{figure}[!htbp]
    \centering
    \includegraphics[width=\textwidth]{SI_AuPd_2panel_side_colors_all_species.pdf}
    \caption{Pourbaix diagrams of AuPd. (a) No added ligands; (b) \ce{NH3}, Gly, and \ce{CN-}. The metal-ion activity is $10^{-4}$. Ligand-containing panels use $[\ce{NH3}]_{\mathrm{tot}}=0.02\,\mathrm{M}$ and $[\mathrm{Gly}]_{\mathrm{tot}}=0.1\,\mathrm{M}$; cyanide-containing panels additionally use $[\ce{CN-}]_{\mathrm{tot}}=10^{-4}\,\mathrm{M}$. The black outline marks pH 10.5--13.5 and $-2.0$ to $2.3$ V vs RHE (converted to SHE and clipped to the plotted range). The side legends list all stable species combinations, including solids, aqueous ions, and metal--ligand complexes. The dashed blue and red lines mark the water reduction and oxidation boundaries, respectively.}
    \label{fig:Au_Pd_Pourbaix_SI}
    \phantomsubcaption\label{fig:AuPd_Pourbaix_NH3_Gly_SI}
    \phantomsubcaption\label{fig:AuPd_Pourbaix_NH3_Gly_CN_SI}
\end{figure}







% % Ti-Ag
% \begin{figure}[!htbp]
%     \centering
%     \begin{subfigure}[b]{0.45\textwidth}
%         \subcaption{}
%         \includegraphics[width=\textwidth]{Figures/pourbaix_diagrams/Ti_Ag/Ti0.5 Ag0.5_NH3=0.02M_Gly=0.1M_CN=0M_activity=1e-04M.png}
%         \label{fig:TiAg_Pourbaix_NH3_Gly_SI}
%        % \par\medskip
%     \end{subfigure}
%     \begin{subfigure}[b]{0.45\textwidth}
%         \subcaption{}
%         \includegraphics[width=\textwidth]{Figures/pourbaix_diagrams/Ti_Ag/Ti0.5 Ag0.5_NH3=0.02M_Gly=0.1M_CN=0.0001M_activity=1e-04M.png}
%         \label{fig:TiAg_Pourbaix_NH3_Gly_CN_SI}
%        % \par\medskip   
%     \end{subfigure}
%     \caption{Pourbaix diagrams of TiAg. (a) $\text{Ion activity}=\num{1e-4}$, $[\ce{NH3}]_{initial}= 0.02$M, $[\text{Gly}]_{initial}=0.1$M,  $[\ce{CN-}]_{initial}=0M$, (b) $\text{Ion activity}=\num{1e-4}$, $[\ce{NH3}]_{initial}= 0.02$M, $[\text{Gly}]_{initial}=0.1$M,  $[\ce{CN-}]_{initial}=\num{1e-4}$M. The black box indicates the experimental condition at applied potential vs RHE = -2 to 2V, pH = 11.5 to 13.5.}
%     \label{fig:Ti_Ag_Pourbaix_SI}
% \end{figure}



%Pd alloys
